\documentclass[convert={svg,command=\unexpanded{{inkscape\space --pdf-poppler\space --export-type="svg"\space --pages=1\space -o\space \outfile\space \infile}}}]{standalone}% 'crop' is the default for v1.0, before it was 'preview'

\usepackage{tikz}
\usepackage{pgfplots}
\usepgfplotslibrary{fillbetween}
\usetikzlibrary{patterns, patterns.meta, positioning, math}

\pgfplotsset{compat=1.18}

\usepackage{tikz-3dplot}
\usetikzlibrary{3d}

\pgfplotsset{
	twoDaxis/.style = {
			no marks,
			samples = 100,
			axis lines = center,
			xlabel = {$x$},
			ylabel = {$y$},
			every axis x label/.style={
					at={(ticklabel* cs:1.05)},
					anchor=west,
				},
			every axis y label/.style={
					at={(ticklabel* cs:1.05)},
					anchor=south,
				},
		},
	threeDaxis/.style = {
			axis z line = none,
			axis lines=center, % remove this and you get weird box behavior, not good for calc 2 students
			xlabel=$x$, % you know this one
			every axis x label/.style={
					at={(ticklabel* cs:1.05)},
					anchor=west,
				}, % place the xlabel in a better spot
			ylabel=$y$, % you know this one
			every axis y label/.style={
					at={(ticklabel* cs:1.05)},
					anchor=north,
				}, % places y label in a better spot
		},
	threeDaxiswithyzswap/.style = {
			axis y line = none,
			axis lines=center, % remove this and you get weird box behavior, not good for calc 2 students
			xlabel=$x$, % you know this one
			every axis x label/.style={
					at={(ticklabel* cs:1.05)},
					anchor=west,
				}, % place the xlabel in a better spot
			zlabel=$y$, % you know this one
			every axis z label/.style={
					at={(ticklabel* cs:1.05)},
					anchor=south, % which side of the label gets pinned to the point 5% past the axis end point?
				}, % places y label in a better spot
			xticklabel style = {font=\tiny, yshift=1ex},
			zticklabel style = {font=\tiny, xshift=1ex},
		}
}



%%%%%%%%%%%%%%%%%%%%%%%%%%%%%%%%%%%

\begin{document}
\begin{tikzpicture}[
		area plot/.style={
				fill opacity=0.4,
				draw=black,
				fill=orange,
				mark=none,
				draw opacity = 0.15,
			}, % sets up the orange slice fill in stuff
	]
	\begin{axis}[
			threeDaxis,
			view={37}{45}, % first number is how much to rotate around z axis for the view, second is how far above/below you are in angle around x (negative is looking up, positive is looking down)
			xticklabel style = {font=\tiny, yshift=0.7ex},
			xtick={0.01,0.523598775598,1.0471975512,1.57079632679,2.09439510239,2.61799387799,3.47640166667,4},
			% need the s here to give a list of labels to use
			xticklabels={$a=x_0$, $x_1$, $x_2$, $x_3$, $\ldots$, $\ldots$, $x_{n-1}$, $x_n$=b},
			ytick=\empty,
			xmin = 0,
			xmax = 5,
			ymin = 0,
			ymax =  5,
			zmin = 0,
			zmax = 3,
		]
		%plot the curve in the x-y plane, keep "samples y=0" or else, and then just plot coords (x,y,z) like parametric curve sorta where x is the parameter. 
		\addplot3 [black,domain=0:4,samples = 100, samples y=0,thick] (x,{0.7*cos(x r)+2},0) node [near end, right] {$f(x)$};
		\addplot3 [black,domain=0:4,samples = 100, samples y=0] (x,{0.7*cos(x r)+2},{0.7*cos(x r)+2});

		% top
		\fill[area plot] (3*3.14159/6,    0,47/20) -- (2*3.14159/6,    0,47/20) -- (2*3.14159/6,47/20,47/20) -- (3*3.14159/6,47/20,47/20) -- cycle;
		% bottom
		\fill[area plot] (3*3.14159/6,    0,    0) -- (2*3.14159/6,    0,    0) -- (2*3.14159/6,47/20,    0) -- (3*3.14159/6,47/20,    0) -- cycle;
		% right
		\fill[area plot] (3*3.14159/6,    0,47/20) -- (3*3.14159/6,    0,    0) -- (3*3.14159/6,47/20,    0) -- (3*3.14159/6,47/20,47/20) -- cycle;
		% left
		\fill[area plot] (2*3.14159/6,    0,47/20) -- (2*3.14159/6,    0,    0) -- (2*3.14159/6,47/20,    0) -- (2*3.14159/6,47/20,47/20) -- cycle;
		% back
		\fill[area plot] (3*3.14159/6,47/20,47/20) -- (2*3.14159/6,47/20,47/20) -- (2*3.14159/6,47/20,    0) -- (3*3.14159/6,47/20,    0) -- cycle;


		\addplot3 [black,domain=0:4,samples = 100, samples y=0] (x,0,{0.7*cos(x r)+2});

		% front
		\fill[area plot] (3*3.14159/6,    0,47/20) -- (2*3.14159/6,    0,47/20) -- (2*3.14159/6,    0,    0) -- (3*3.14159/6,    0,    0) -- cycle;

		\node [circle,fill,inner sep=0.7pt] at (3.14159/3,0,47/20) {};
		\node [circle,fill,inner sep=0.7pt] at (3.14159/3,47/20,47/20) {};

		\tikzmath {
			% set the number of rectangles (it's off by 1 but I don't care)
			\N = 1;
		}

		%set up the loop
		\pgfplotsinvokeforeach{0,...,\N} {
			\tikzmath{
				% calculate the current x value for each iteration (otherwise it acts weird) (the 4 is hard coded the width of theinterval and "#1" is like "i" in a C loop, it is "this iteration number")
				\currentx{#1} = 4/\N*#1;
				% cacluctate the current y value for each iteration (we use it multiple times so only calculate it once here)
				\currenty{#1} = 0.7*cos(\currentx{#1} r)+2;
			}
			% plot the slice for this iteration by plotting the 4 coordinates of the square's vertices and use closedcycle to connect them, using the area plot style we set up earlier to color them orange and all that
			\addplot3 [domain=0:4,samples = 100, samples y=0,black] coordinates {(\currentx{#1},0,0)  (\currentx{#1},\currenty{#1},0)  (\currentx{#1},\currenty{#1},\currenty{#1})  (\currentx{#1},0,\currenty{#1})} \closedcycle;
		}

	\end{axis}
\end{tikzpicture}
\end{document}
